\documentclass[11pt,a4paper]{article}
\usepackage[utf8]{inputenc}
\usepackage[spanish,es-noshorthands]{babel}
\usepackage{amsmath}
\usepackage{amsfonts}
\usepackage{amssymb}
\usepackage{graphicx}
\usepackage[hidelinks]{hyperref}
\usepackage{listings}
\usepackage{color}
\usepackage{geometry}
\geometry{margin=1in}

\definecolor{lightgray}{rgb}{0.95, 0.95, 0.95}
\definecolor{codegray}{rgb}{0.5, 0.5, 0.5}
\definecolor{codeblue}{rgb}{0.1, 0.2, 0.6}

\lstset{
    backgroundcolor=\color{lightgray},
    basicstyle=\small\ttfamily,
    breakatwhitespace=false,
    breaklines=true,
    captionpos=b,
    commentstyle=\color{codegray},
    keepspaces=true,
    keywordstyle=\color{codeblue},
    numbers=left,
    numbersep=5pt,
    numberstyle=\tiny\color{codegray},
    showspaces=false,
    showstringspaces=false,
    showtabs=false,
    tabsize=2,
    frame=single,
    rulecolor=\color{lightgray}
}

\title{Manual de Operación y Despliegue \\ \large Cloudflare Serverless Agent}
\author{Espacio de Trabajo - Electrónica}
\date{\today}

\begin{document}

\maketitle

\tableofcontents
\newpage

\section{Introducción}
Este manual describe el flujo de trabajo para el desarrollo, prueba y despliegue del \textbf{Cloudflare Serverless Agent} en la dirección de producción \href{https://cloudflare-agent.cero-agents.workers.dev/}{cloudflare-agent.cero-agents.workers.dev}. El proyecto local está ubicado en la ruta:
\begin{verbatim}
/home/cero/MEGA/VS_CODE_WORKSPACE/ELECTRONICA/Proyectos/cloudflare-agent
\end{verbatim}

\section{Estructura del Proyecto}
El proyecto consta de los siguientes archivos principales:
\begin{itemize}
    \item \textbf{package.json}: Define los scripts de npm y las dependencias (como \texttt{wrangler} y \texttt{pdf-lib}).
    \item \textbf{wrangler.toml}: Configuración de Wrangler para el Worker de Cloudflare, especificando el nombre del servicio (\texttt{cloudflare-agent}) y el punto de entrada (\texttt{src/index.ts}).
    \item \textbf{src/index.ts}: Código fuente TypeScript principal que implementa los endpoints (\texttt{/}, \texttt{/status}, \texttt{/chat}, \texttt{/quote}).
    \item \textbf{.dev.vars}: Archivo local (excluido de git) donde se configuran los tokens/claves de API para el entorno de desarrollo local.
\end{itemize}

\section{Flujo de Desarrollo Local}
Antes de realizar un despliegue en producción, es recomendable probar y validar los cambios de forma local.

\subsection{Inicio del Servidor de Desarrollo}
Para arrancar el entorno de desarrollo local, ejecuta desde la carpeta raíz del proyecto:
\begin{lstlisting}[language=bash]
npm run dev
\end{lstlisting}
O de forma directa mediante Wrangler:
\begin{lstlisting}[language=bash]
npx wrangler dev
\end{lstlisting}

Esto iniciará un servidor local (por lo general en \texttt{http://localhost:8787}) que emulará el entorno del Worker de Cloudflare. Este servidor leerá automáticamente las variables de entorno definidas en el archivo local \texttt{.dev.vars}.

\section{Flujo de Despliegue a Producción}
Cuando los cambios estén validados y listos para producción, sigue el siguiente procedimiento:

\subsection{Autenticación (Si es requerida)}
Si no has iniciado sesión o si la sesión con Cloudflare ha expirado, debes autenticarte en Wrangler ejecutando:
\begin{lstlisting}[language=bash]
npx wrangler login
\end{lstlisting}
Esto abrirá el navegador para ingresar tus credenciales y autorizar la CLI en tu máquina local.

\subsection{Despliegue del Código}
Para compilar y subir los cambios a la infraestructura de Cloudflare, ejecuta:
\begin{lstlisting}[language=bash]
npm run deploy
\end{lstlisting}
O bien:
\begin{lstlisting}[language=bash]
npx wrangler deploy
\end{lstlisting}
Una vez completado de manera exitosa, los cambios se verán reflejados inmediatamente en \href{https://cloudflare-agent.cero-agents.workers.dev/}{https://cloudflare-agent.cero-agents.workers.dev/}.

\section{Configuración de Secrets en Producción}
El código en \texttt{src/index.ts} depende de tres variables de entorno fundamentales para realizar llamadas a los LLMs:
\begin{itemize}
    \item \texttt{GROQ\_API\_KEY}
    \item \texttt{OPENROUTER\_API\_KEY}
    \item \texttt{GOOGLE\_API\_KEY}
\end{itemize}

Dado que \texttt{.dev.vars} no se despliega en producción por seguridad, debes configurar estas claves de API en Cloudflare de la siguiente manera:

\begin{lstlisting}[language=bash]
npx wrangler secret put GROQ_API_KEY
npx wrangler secret put OPENROUTER_API_KEY
npx wrangler secret put GOOGLE_API_KEY
\end{lstlisting}

Ejecuta cada comando individualmente e ingresa el valor de la clave correspondiente cuando la CLI te lo solicite.

\end{document}
